\begin{frame}{行情计算方法}\relax
\Intro{采用腾讯财经前复权日线，对照同日期的指数描述市场表现。}
\begin{columns}[T,onlytextwidth]
\column{.48\textwidth}
\Head{交易窗口}
披露时间只有日期时，采用之后首个交易日作为 T$_0$。基期为 T$_{-1}$ 收盘。
\Gap
1、3、5 日窗口分别截至 T$_0$、T$_{+2}$、T$_{+4}$；按各市场实际交易日顺延。
\Gap
GLM-5：T$_0$ 为 2026-02-13，5 日窗口截至 2026-02-24。
\column{.48\textwidth}
\Head{累计收益与基准}
\[
  R_k=\frac{P_{T_{k-1}}}{P_{T_{-1}}}-1
\]
基准使用同一组日期；相对基准值为累计收益的简单百分点差。
\Gap
港股以恒生指数、A 股以\mbox{科创50}为基准；未上市或行情不足时不计算收益。
\end{columns}
\end{frame}
